\documentclass[aspectratio=169,10pt,spanish]{beamer}

\input{../../../_preambulo_beamer.tex}
\input{../../../_comandos_beamer.tex}

\selectlanguage{spanish}

\title{MA1001B - Modelación Estadística para la Toma de Decisiones}
\subtitle{Lección 48: Inferencia de regresión, residuos y decisiones}
\author{}
\institute{Tecnol\'ogico de Monterrey}
\date{}

\begin{document}

\begin{frame}[plain]
  \vspace{1.2cm}
  {\Large\bfseries MA1001B - Modelación Estadística para la Toma de Decisiones\par}
  \vspace{0.7cm}
  {\large Lección 48: Inferencia de regresión, residuos y decisiones\par}
  \vspace{0.7cm}
  {\color{TecRojo}\rule{\textwidth}{1.2pt}}
  \vspace{0.8cm}

  Facultad MA1001B\\[0.3cm]
  Periodo\\[0.3cm]
  Tecnol\'ogico de Monterrey
\end{frame}

\begin{frame}{La pregunta de inferencia}
  \begin{alertblock}{Pregunta guía}
    ¿Es una pendiente ajustada de horas extra distinta de 0, y puede un
    fin de semana de emergencia crear esa significancia por sí solo?
  \end{alertblock}

  \vspace{0.6em}
  Al final de esta lección, usted puede:
  \begin{itemize}
    \item calcular $EE(b_1)$ y la prueba $t$ de $H_0:\beta_1=0$;
    \item leer una gráfica residuo-frente-a-ajustado;
    \item reajustar tras eliminar una semana influyente;
    \item rechazar tratar un fin de semana como política de semanas
      ordinarias.
  \end{itemize}

  \vfill
  \scriptsize Todas las horas extra y conteos de unidades usados hoy son sintéticos. No se usan datos reales de empresa.
\end{frame}

\begin{frame}{Treinta y nueve semanas ordinarias y una emergencia}
  \small
  \begin{center}
    \begin{tabular}{lr}
      \toprule
      \textbf{Pieza} & \textbf{En esta muestra} \\
      \midrule
      Semanas ordinarias & 39, horas extra en $[6,\ 22]$ \\
      Fin de semana de emergencia & $(x,y)=(48,\ 42)$ \\
      $H_0$ & $\beta_1=0$ \\
      $H_1$ & $\beta_1\neq 0$ \\
      \bottomrule
    \end{tabular}
  \end{center}

  \vspace{0.5em}
  \begin{exampleblock}{Por qué importa el fin de semana}
    Un $x$ de alto apalancamiento puede inclinar la recta de mínimos
    cuadrados. La inferencia en las 40 semanas no es automáticamente
    inferencia sobre operaciones ordinarias.
  \end{exampleblock}
\end{frame}

\begin{frame}[fragile]{Paso 1: Pendiente, EE y $t$}
  \begin{columns}[T]
    \begin{column}{0.40\textwidth}
      \small
      \begin{lstlisting}
t = slope / se_slope
p = two_sided_t(t, df=38)
      \end{lstlisting}

      \begin{block}{Salida verificada}
        $b_1=0.434381$\\
        $EE=0.072494$, $t=5.991932$\\
        $p=5.839701\times10^{-7}$
      \end{block}
    \end{column}
    \begin{column}{0.60\textwidth}
      \centering
      \includegraphics[width=\textwidth]{../figuras/inferencia_regresion_01_pendiente_ee_t.png}
      \vspace{0.2em}
      \scriptsize Recta ajustada $n=40$ del Paso 1
    \end{column}
  \end{columns}
\end{frame}

\begin{frame}{Los residuos localizan la semana inusual}
  \small
  La media residual es $0.000000$ por la identidad de mínimos cuadrados.
  La DE residual es $3.140288$. El residuo del fin de semana de emergencia
  es $8.570594$, el residuo absoluto más grande de la muestra.
  \[
    e_i=y_i-\hat y_i=y_i-(12.579109+0.434381\,x_i).
  \]
  Una gráfica de residuos es una herramienta de decisión: puede mostrar un
  abanico, una curva o una semana que no pertenece con las demás.
\end{frame}

\begin{frame}[fragile]{Paso 2: Residuo frente a ajustado}
  \begin{columns}[T]
    \begin{column}{0.40\textwidth}
      \small
      \begin{lstlisting}
e = y - (b0 + b1 * x)
      \end{lstlisting}

      \begin{block}{Salida verificada}
        Media residual $0.000000$\\
        Emergencia $e=8.570594$\\
        El $|e|$ más grande es esa semana
      \end{block}
    \end{column}
    \begin{column}{0.60\textwidth}
      \centering
      \includegraphics[width=\textwidth]{../figuras/inferencia_regresion_02_grafica_residuos.png}
      \vspace{0.2em}
      \scriptsize Gráfica de residuos del Paso 2
    \end{column}
  \end{columns}
\end{frame}

\begin{frame}{Paso 3: La pendiente significativa desaparece}
  \begin{columns}[T]
    \begin{column}{0.42\textwidth}
      \small
      Sin el fin de semana: $n=39$, $b_1=0.088445$, $t=1.070131$,
      $p=0.291492$, $R^2=0.030022$.

      \vspace{0.4em}
      No rechazar $H_0:\beta_1=0$.

      \vspace{0.5em}
      \textbf{Límite:} un punto influyente puede crear una pendiente
      significativa que desaparece al quitarlo. La prueba $n=40$ no es una
      decisión sobre semanas extra ordinarias.
    \end{column}
    \begin{column}{0.58\textwidth}
      \centering
      \includegraphics[width=\textwidth]{../figuras/inferencia_regresion_03_punto_influyente.png}
      \vspace{0.2em}
      \scriptsize Rectas $n=40$ frente a $n=39$ del Paso 3
    \end{column}
  \end{columns}
\end{frame}

\begin{frame}{Verificación de conceptos}
  \begin{enumerate}
    \item Escriba el estadístico $t$ para $H_0:\beta_1=0$ en términos de
      $b_1$ y $EE(b_1)$.
    \item ¿Por qué la media residual 0 no es, por sí sola, un aprobado
      diagnóstico?
    \item ¿Qué ocurre con $p$ cuando se elimina el fin de semana de
      emergencia?
    \item ¿Por qué la pendiente $n=40$ es el número equivocado de política
      de semanas ordinarias?
  \end{enumerate}

  \vspace{0.6em}
  \begin{block}{Conexión}
    Esto cierra la secuencia de regresión simple. Los predictores múltiples
    y los experimentos diseñados son cursos posteriores; el límite aquí es
    la influencia.
  \end{block}

  \vfill
  \scriptsize Fuente: OpenStax, \textit{Introductory Business Statistics 2e},
  Secciones 13.2 y 13.6. CC BY-NC-SA.
\end{frame}

\appendix



\end{document}
